\section{The \texttt{Arm} McXtrace Component}
Arm/optical bench

\subsection*{Identification}
\begin{itemize}
  \item \textbf{Author:} Kim Lefmann and Kristian Nielsen
  \item \textbf{Origin:} Risoe
  \item \textbf{Date:} September 2009
\end{itemize}

\subsection*{Description}
An arm does not actually do anything, it is just there to set up a new coordinate system.

Example: Arm()

\subsection*{Input parameters}
Parameters in \textbf{boldface} are required; the others are optional.

\begin{longtable}{p{0.22\textwidth}p{0.12\textwidth}p{0.46\textwidth}p{0.14\textwidth}}
\toprule
\textbf{Name} & \textbf{Unit} & \textbf{Description} & \textbf{Default} \\
\midrule
\endhead
\bottomrule
\end{longtable}

\subsection*{Links}
\begin{itemize}
  \item Component source code found in file \texttt{Arm.comp}.
\end{itemize}
\IfFileExists{optics/Arm_static.tex}{\input{optics/Arm_static.tex}}{}